\chapter{Victimology and Identity}
\label{chap:victimology-and-identity}

\classification{Evidence status: victimology and identity assessment. Confidence baseline: High for post-discovery observations, indexed possessions, movement records, and forensic identifiers where source-mapped; Moderate for OCR-dependent details and modern public-source identity research; Unknown where pre-death biography, family, occupation, residence, and official identity status are not established in the repository. The Carl Charles Webb attribution is treated as a researcher/lab-led identification lead pending official coroner or South Australia Police confirmation.}

\section{Victim Overview}

The demonstrable subject of this chapter is the unknown man found at Somerton. In the acquired 1949 inquest record, he remains unidentified. The project labels Somerton Man and Tamam Shud man are research labels, not original legal identity.

The repository supports a strong post-discovery victimology profile and a weak pre-death biography. It can document physical description, body condition, clothing, possessions, tickets, witness observations, fingerprints, dental charting, body photographs, death-mask or bust creation, and biological specimens. It cannot yet verify date of birth, place of birth, family, residence, employment, education, military service, finances, or associates for the deceased from acquired primary records.

The current modern identity position is controlled by the identity-status guidance: a serious researcher/lab-led Carl Charles Webb attribution exists, strengthened by Mission 11 OSINT, but the repository has not acquired a final coroner, South Australia Police, or Forensic Science SA report closing the official identity question.

\begin{table}[htbp]
\centering
\caption{Victim profile summary}
\label{tab:victim-profile-summary}
\begin{tabularx}{\textwidth}{>{\raggedright\arraybackslash}p{0.28\textwidth}>{\raggedright\arraybackslash}p{0.28\textwidth}>{\raggedright\arraybackslash}X}
\toprule
Profile area & Current evidence position & Confidence and limitation \\
\midrule
1949 identity & Unknown man found at Somerton. & Verified as unidentified in acquired inquest evidence. \\
Modern identity lead & Carl Charles Webb attribution. & Researcher/lab-led lead; official-status gap remains. \\
Birth and family & Unknown in acquired primary archive. & No verified birth, family, or civil record for deceased. \\
Occupation and skills & Unknown. & Possessions do not establish occupation or technical skill. \\
Physical description & Male, tallish, about mid-40s, greying hair, good physical condition, missing teeth, cared-for nails/feet. & Partially verified; exact measurements and chart transcription incomplete. \\
Possessions & Clothing, tickets, cigarettes/matches, combs, gum, Tamam Shud slip, suitcase/contents, identification records, specimens. & Strong for existence of many entries; custody and exact details vary. \\
Movements & Ticket evidence, beach observations, discovery, transport, postmortem, suitcase lead. & Strong for post-discovery and record events; weak for pre-death route reconstruction. \\
Relationships & Official witnesses and investigators; conditional modern Webb, Thomson, Boxall leads. & Strong for official contacts; weak or conditional for pre-death personal relationships. \\
\bottomrule
\end{tabularx}
\end{table}

\section{Physical Characteristics}

The acquired victimology profile records the deceased as a male in approximately mid-life. Dwyer described the body as tallish, about 45 years, and in good physical condition. The forensic identifiers register cautions that exact height and build are not precisely measured in the extracted medical pages and that an autopsy sheet or police physical description is still needed for exact measurements.

Hair is documented as greying in the victim profile, but the project does not yet have a complete primary body-description sheet. Eye colour is not established in the acquired victimology package. Hands and feet are relevant only at the level of recorded observation: Dwyer described signs of care in fingernails and feet, and the forensic identifiers register records a light scar on the left upper arm, an uncertain vaccination mark, and two recent right-hand abrasions.

Dentition is one of the stronger identifying features. Dwyer prepared a missing-teeth chart, and a separate tooth-chart reproduction is preserved as DOC-0017. Several teeth were missing, and Dwyer stated that the missing teeth would be noticeable if the man laughed but not necessarily if speaking. The limitation is that the chart image has no OCR and needs specialist or manual transcription before precise dental notation is used in prose.

Fingerprint evidence is also strong for process. Durham fingerprinted the deceased and circulated copies to Australian states, New Zealand, and important overseas bureaus. A fingerprint chart reproduction is preserved, but the original chart, complete comparison logs, and bureau responses are not present in the repository.

\section{Lifestyle Indicators}

Lifestyle indicators are treated as observations, not personality conclusions. Clothing was described as fairly good quality, and the body was described as in good physical condition with cared-for nails and feet. Shoes and clothing appear in the physical-evidence and possession registers. These observations do not establish social class, occupation, motive, planning, or personal habits by themselves.

Possessions show a narrow post-discovery inventory rather than a complete life history. The body-property layer includes clothing, transport tickets, cigarette-related items, matches, combs, chewing gum, and the later-located Tamam Shud slip. The suitcase layer includes clothing/textiles, tools or personal articles, pencils, thread, and related property. The repository records these as possession and association evidence, not proof of occupation or ownership in every case.

Travel indicators are limited to ticket evidence and movement records. The Henley Beach railway ticket and Adelaide-Somerton bus ticket are high-confidence property records, but the purchaser, holder, boarding point, alighting point, and full route remain unknown. The suitcase was associated with Adelaide Railway Station luggage-office evidence, but depositor identity is unknown.

Habit evidence is weak. Smoking is suggested by cigarette evidence and possible finger-staining testimony, but the behavioural profile states that habit-level wording requires manual transcription before final use. Grooming and cleanliness are supported as witness or medical observations only.

Financial indicators are also weak. A 2s 6d cash entry appears in OCR-dependent testimony, but it requires visual confirmation before being used as a firm detail. The repository does not establish financial status.

\section{Movement Reconstruction}

The movement reconstruction contains 11 records. The first two are ticket records from 30 November 1948: the Henley Beach rail ticket issue window during Townsend's 6:15 a.m. to about 2:00 p.m. shift, and the bus ticket associated with the Adelaide-Somerton route departing the Railway Station at about 11:15 a.m. Both records are high confidence as ticket events, but neither identifies the purchaser or holder.

The next movement entries concern the interval before the evening beach sightings. The coroner recorded that no one came forward to say the man was seen at Somerton between the bus-arrival interval and 7:00 p.m. This is negative evidence only. It depends on the completeness of the witness canvass and does not prove absence.

The evening records are witness observations. At about 7:00 p.m., a man was seen near the steps or seawall and was reported to raise his right arm before it fell. At about 7:20 to 8:00 p.m., Strapps and Neill observed a man at the beach and saw no movement. These entries are high confidence as observations, but identity continuity remains unresolved.

The 1 December entries begin with Bennett's approximate medical time estimate, then the morning report and discovery, then movement of the body by police ambulance to Royal Adelaide Hospital and City Mortuary. The movement layer also records Dwyer's 2 December postmortem and the January 1949 suitcase event at Adelaide Railway Station. These are documentary and investigative anchors, not a complete pre-death route.

\begingroup
\scriptsize
\setlength{\tabcolsep}{2pt}
\renewcommand{\arraystretch}{1.12}
\begin{longtable}{>{\raggedright\arraybackslash}p{0.08\textwidth}>{\raggedright\arraybackslash}p{0.12\textwidth}>{\raggedright\arraybackslash}p{0.14\textwidth}>{\raggedright\arraybackslash}p{0.27\textwidth}>{\raggedright\arraybackslash}p{0.10\textwidth}>{\raggedright\arraybackslash}p{0.17\textwidth}}
\caption{Movement reconstruction}\label{tab:chapter8-movement-reconstruction}\\
\toprule
ID & Date & Time & Movement or event & Confidence & Principal uncertainty \\
\midrule
\endfirsthead
\toprule
ID & Date & Time & Movement or event & Confidence & Principal uncertainty \\
\midrule
\endhead
MOV-0001 & 1948-11-30 & 06:15--14:00 & Henley Beach rail ticket issued during Townsend's shift window. & High & Purchaser and exact issue time unknown. \\
MOV-0002 & 1948-11-30 & 11:15 a.m. route & Bus ticket associated with body issued on Adelaide-Somerton route. & High & Holder, boarding, and alighting unknown. \\
MOV-0003 & 1948-11-30 & Before 7:00 p.m. & No witness came forward for interval between bus arrival and 7:00 p.m. & Moderate & Negative evidence only. \\
MOV-0004 & 1948-11-30 & 7:00 p.m. & Man seen near steps or seawall raising right arm. & High & Face not seen; identity continuity unresolved. \\
MOV-0005 & 1948-11-30 & About 7:20--8:00 p.m. & Strapps and Neill observed a man at the beach with no movement. & High & Limited evening observation. \\
MOV-0006 & 1948-12-01 & About 2:00 a.m. estimate & Bennett time-of-death estimate. & Moderate & Medical estimate; place unknown. \\
MOV-0007 & 1948-12-01 & About 6:35--6:45 a.m. & Morning witness activity and Moss response. & High & Exact reporting sequence needs confirmation. \\
MOV-0008 & 1948-12-01 & About 7:00 a.m. & Body found on shore at Somerton. & Very High & First-discovery details need visual confirmation. \\
MOV-0009 & 1948-12-01 & Exact time unknown & Body conveyed to Royal Adelaide Hospital and City Mortuary. & High & Exact transfer times not indexed. \\
MOV-0010 & 1948-12-02 & 7:30 a.m. & Dwyer performed postmortem at City Mortuary. & High & Medical detail needs manual transcription. \\
MOV-0011 & 1949-01-14/19 & Unknown & Leane saw and later took possession of unclaimed suitcase. & High & Post-police custody unresolved. \\
\bottomrule
\end{longtable}
\endgroup

\section{Relationships}

The relationship map contains 18 entries. Fourteen are official, witness, transport, medical, or investigative relationships created by the case process. These are not personal relationships. They include the beach witnesses, Moss, Hall, Holdernesse, Strapps, Neill, North, Leane, Durham, Townsend, Dwyer, Cowan, Lawson, and Brown.

No verified family, friend, employer, or pre-death associate relationship is established in the current acquired primary-source archive. The relationship map records family, friends, and employer as unknown.

Jessica Thomson is retained only as a partially verified or unable-to-verify later-source connection through reported phone or Rubaiyat material. The repository does not contain an acquired primary interview or residence source sufficient to establish a personal relationship with the deceased. Robin Thomson is not used here as an established relationship to the deceased unless future repository evidence supports that link.

Alf Boxall is retained as a later Rubaiyat-related lead, not as a verified direct relationship to the deceased. The acquired Boxall military record does not establish a direct relationship to the unknown man. The Boxall/Jestyn material belongs in documentary and intelligence-assessment contexts unless stronger primary relationship evidence is acquired.

The Carl Charles Webb attribution appears in the relationship map as a proposed modern identity, not a relationship. If official confirmation is later acquired, Webb-related family, marriage, divorce, and residence records would become direct biographical material for the victimology chapter. Until then, those records remain conditional identity leads.

\begin{table}[htbp]
\centering
\caption{Relationship map summary}
\label{tab:relationship-map-summary}
\begin{tabularx}{\textwidth}{>{\raggedright\arraybackslash}p{0.28\textwidth}r>{\raggedright\arraybackslash}X}
\toprule
Relationship class & Count & Treatment in this chapter \\
\midrule
Official or witness contact after death & 14 & Documented case-process relationships, not personal relationships. \\
Modern identity attribution & 1 & Webb attribution pending official confirmation. \\
Later reported Thomson/Rubaiyat lead & 1 & Partially verified or unable to verify; not a demonstrated relationship. \\
Later Boxall/Rubaiyat lead & 1 & Unable to verify as direct relationship to deceased. \\
Family, friends, employer & 1 & Unknown in acquired primary archive. \\
\bottomrule
\end{tabularx}
\end{table}

\section{Identity Assessment}

Historic identification attempts are documented through fingerprints, photographs, dental charting, public circulation, and missing-person inquiry, but they did not identify the man in the 1949 inquest record. Durham's testimony supports fingerprinting, photography, and circulation. Dwyer's and the dental-chart records support dental identification material. The death certificate and inquest record preserve unidentified status in the acquired archive.

The modern identity evidence is materially stronger than ordinary media speculation. Mission 11 records that Abbott announced in July 2022 that the man was Carl Charles Webb, based on DNA from hairs in the plaster bust and genealogical matching. It also records Astrea Forensics' account of extracting, sequencing, and genotyping DNA from a single 5 cm rootless hair strand, and Fitzpatrick's description of GEDmatch and family-tree reconstruction with Keane-line and maternal-side matching.

Mission 11 also records biographical leads associated with Webb: ABC-reported civil-record details for Carl Webb, including birth on 16 November 1905 in Footscray, marriage to Dorothy Jean Robertson on 4 October 1941, and occupation as instrument maker; and a 5 October 1951 Trove legal notice in which Dorothy Jean Webb instituted divorce proceedings against Carl Webb, formerly of Bromby Street, South Yarra, then of parts unknown.

These Webb records are important, but the chapter must preserve their evidentiary status. The repository does not contain the final coroner, South Australia Police, or Forensic Science SA determination. As of the Mission 11 OSINT update, ABC reported on 6 July 2026 that South Australia Police had not yet presented its findings to the coroner. No peer-reviewed public DNA identification paper or official coroner finding was located in that OSINT pass.

The controlled wording is therefore: researcher/lab-led Carl Charles Webb attribution, pending official coroner or South Australia Police determination. The chapter avoids closed-identity language and does not convert public DNA reporting into legal identification.

\begin{table}[htbp]
\centering
\caption{Identity attempts and status timeline}
\label{tab:identity-attempts-timeline}
\begin{tabularx}{\textwidth}{>{\raggedright\arraybackslash}p{0.16\textwidth}>{\raggedright\arraybackslash}p{0.30\textwidth}>{\raggedright\arraybackslash}p{0.22\textwidth}>{\raggedright\arraybackslash}X}
\toprule
Date or period & Identity action or lead & Evidentiary status & Limitation \\
\midrule
1948--1949 & Photographs, fingerprints, dental charting, public circulation, and missing-person inquiries. & Documented in acquired inquest/index registers. & No identification in acquired 1949 record. \\
1949 & Rubaiyat and Boxall/Jestyn leads. & Later-source/documentary leads. & No acquired primary proof of direct relationship to deceased. \\
2021 & Official exhumation for modern forensic/DNA work. & Official action in OSINT/source register. & Full official forensic report not acquired. \\
2022 & Abbott/Fitzpatrick/Astrea Webb attribution. & Serious researcher/lab-led identification lead. & Public method details incomplete; official legal determination not acquired. \\
2025--2026 & High-quality reporting and ABC current-status update. & Supports public status and official-status gap. & No final coroner/SAPOL/FSSA report in repository. \\
\bottomrule
\end{tabularx}
\end{table}

\section{Victimology Assessment}

What is known with the strongest support is post-discovery. The deceased was an unidentified male found at Somerton Beach, with documented clothing and property, ticket evidence, witness observations, fingerprinting, dental charting, body photography, death-mask or bust creation, and biological specimens submitted for testing.

Supported inferences are limited. The transport tickets support movement leads but not a complete route. The suitcase supports an association lead but not complete ownership or depositor identity. Clothing quality, cleanliness, and grooming observations support appearance description but not social class, personality, or motive. The Tamam Shud slip placement is documented; intent is unknown.

Modern identity research now provides a strong public Webb lead, but not a project-verified legal identity. If official confirmation is acquired, Webb-related civil and family records will become central to this chapter. Until then, they remain conditional identity material.

The largest unknowns are pre-death biography, official identity status, exact pre-death movements, personal relationships, residence, employment, finances, purpose for being at Somerton, and whether the man seen on the evening of 30 November was definitively the deceased.

\begin{table}[htbp]
\centering
\caption{Victimology confidence matrix}
\label{tab:victimology-confidence-matrix}
\begin{tabularx}{\textwidth}{>{\raggedright\arraybackslash}p{0.30\textwidth}>{\raggedright\arraybackslash}p{0.20\textwidth}>{\raggedright\arraybackslash}X}
\toprule
Assessment area & Confidence & Basis and limitation \\
\midrule
Unidentified status in 1949 record & High & Inquest and death-certificate layer retain unknown identity. \\
Physical description & Moderate & Supported by testimony and identifiers; exact measurements incomplete. \\
Possession inventory & Moderate--High & Many entries source-mapped; custody and OCR detail vary. \\
Movement reconstruction & Moderate--High & Strong post-discovery anchors; pre-death route incomplete. \\
Relationships & Low for personal biography & Official contacts verified; personal contacts unknown or conditional. \\
Webb attribution & Moderate-High as research lead & Strong public DNA/genealogy pathway; official determination not acquired. \\
Pre-death occupation, residence, family, finances & Unknown & No acquired primary records verify these for the deceased. \\
\bottomrule
\end{tabularx}
\end{table}
